% === Begin Label Data ===
% ID: 816
% 难度星级: 
% 标签: 
% 备注: 
% 组卷引用次数: 0
% === End  Label Data ===

\begin{problem}{2018}{G}{全国III卷（文）}{6}{三角函数}
函数 $ f(x)=\dfrac{\tan x}{1+\tan^2x} $ 的最小正周期为 (\hspace{1cm})
\begin{choices}
\choice{{$\dfrac{\pi}{4}$}}
\choice{{$\dfrac{\pi}{2}$}}
\choice{{$\pi$}}
\choice{{$2\pi$}}
\end{choices}
\end{problem}

\begin{answer}
C
\end{answer}

\begin{solutions}
 思路1：直接从函数的表达式不易找到函数的最小正周期，因此需按照“化弦”的基本思想，将函数改为由正弦、余弦表示，进一步化为同角、同名的三角函数并求出最小正周期．  
因 $ \tan x=\dfrac{\sin x}{\cos x} $，由题设可得  
$$
f(x)=\dfrac{\tan x}{1+\tan^2x}=\dfrac{\dfrac{\sin x}{\cos x}}{1+\dfrac{\sin^2x}{\cos^2x}}=\dfrac{\sin x\cos x}{\sin^2x+\cos^2x}=\sin x\cos x=\dfrac{1}{2}\sin 2x.
$$  
因为 $ \sin x $ 的最小正周期为 $ 2\pi $，故 $ f(x)=\dfrac{1}{2}\sin 2x $ 的最小正周期为 $ \pi $．因此选项 C 正确．  

思路2：由题设知函数 $ f(x) $ 的定义域为 $ \{x\mid x\in\mathbb{R} \text{ 且 } x\ne k\pi+\dfrac{\pi}{2},k\in\mathbb{Z}\} $．若最小正周期为 $ \dfrac{\pi}{2} $，即有 $ f(x+\dfrac{\pi}{2})=f(x) $，令 $ x=0 $，可得 $ f(0)=f(\dfrac{\pi}{2}) $，但 $ f(\dfrac{\pi}{2}) $ 无意义，因此选项 B 不正确．  
同理，若 $ f(x+\dfrac{\pi}{4})=f(x) $，令 $ x=\dfrac{\pi}{4} $，可得 $ f(\dfrac{\pi}{4})=f(\dfrac{\pi}{2}) $，因此选项 A 也不正确．  
又 $ \pi $ 和 $ 2\pi $ 均是 $ \tan x $ 和 $ \tan^2x $ 的周期，即均有 $ f(x+\pi)=f(x) $ 与 $ f(x+2\pi)=f(x) $，但题设需求出函数的最小正周期，因此选项 C 正确．
\end{solutions}